% ===== §9 =====
\section{The Two Faces of the Real}
\label{sec:twofaces}

\noindent
This section establishes that the gap at the source of the six limits has two faces, one genetic and one logical, and that they are the same gap. The aim is to show that the pre-symbolic remainder (a part cut off at the origin) and the impossibility of complete symbolization (a closure that can never be attained) are not two conditions but one, so that the single name \emph{manque} covers both. The section takes the genetic face first, draws the logical face out of it as a consequence rather than a separate posit, reviews the reading that would keep them apart, and states the claim.

\subsection*{The genetic face}

Genetically, as the previous section set out, a remainder is cut off in the constitution of the speaking subject. Something is left outside the signifier when the subject is taken up into the symbolic order. This is a fact about how a speaking subject comes to be: it comes to be by a cut, and the cut leaves a portion in the Real. The genetic face is a claim about origin. There was, in the constitution of the subject, an operation of separation, and that operation has a remainder.

\subsection*{The logical face, drawn from the genetic}

From this the logical face follows, and the order matters: the logical is not an independent thesis but the consequence of the genetic. Because a remainder is left outside the signifier in the very operation that installs the symbolic order, the complete closure of that order is not a task still to be finished but a thing that cannot be. The symbolic cannot close over the whole, because its own installation was the leaving-out of a part. The impossibility is not empirical, not a matter of the signifier having so far failed to cover everything and perhaps succeeding later with better instruments. It is structural. The operation that builds the symbolic order is the same operation that produces what the order cannot contain, so the order runs on a constitutive non-closure. Fink draws the consequence exactly: there is no metalanguage, no signifying system able to close upon itself and speak its own whole, because every such system is founded on the exclusion that founds it \citep{fink1995subject}. This is the Real in its logical sense, what Lacan in the late teaching approached as the impossible, that which does not cease not to be written \citep{lacan1998xx}: not the merely not-yet-said but the impossible-to-say-completely, the point at which symbolization meets its own limit as a limit and not as a frontier.

\subsection*{Against keeping the two faces apart}

It might seem cleaner to keep the two faces separate, to treat the pre-symbolic remainder as a developmental posit, a something-before that could in principle be dispensed with, and to rest the impossibility of closure on a purely logical argument about self-reference, in the manner of the formal results on systems that cannot prove their own consistency. Some readings of Lacan lean this way, treating the Real as either the mythical before or the logical limit, and choosing between them \citep{copjec1994read,chiesa2007subjectivity}. The present paper declines the choice, and for a reason internal to the account. The logical impossibility, taken alone, is a fact about formal systems; it does not by itself explain why there should be a \emph{drive}, a felt pressure, at the point of impossibility. The genetic cut, taken alone, is a story about origin; it does not by itself establish that closure is impossible in principle rather than merely unattained. Held together, each supplies what the other lacks. The genetic cut explains why the impossibility is inhabited, why there is a pressure and not merely a formal gap: because the remainder is a part of the subject's own being, left outside and pressing to return. The logical impossibility explains why the cut can never be healed, why the pressure is permanent and not a wound that time might close. The two faces are one because only together do they give the phenomenon the paper is after, an impossibility that is also a drive.

\subsection*{The single gap}

The two faces are one gap. \emph{Manque} (the lack, the constitutive want) names it. It is genetic in that a part was cut off, and it is logical in that, a part having been cut off, closure became impossible. What was left out at the origin is what can never be gathered in at the end, and the reason no sentence in Part One could reach its core is that the core, in every case, was manque: the remainder whose being-cut-off is one with the impossibility of its ever being said.

\begin{claim}[The single gap]
The pre-symbolic remainder and the impossibility of complete symbolization are not two conditions but the two faces of one manque. Because the subject is constituted by a cut that leaves a part in the Real, the symbolic order cannot close; the genetic want grounds the logical impossibility, and the logical impossibility makes the genetic want permanent. Neither face alone yields an impossibility that is also a drive; held together, they do.
\end{claim}
